\documentclass{article}
\usepackage{amsmath,amssymb,amsthm,mathtools}
\newtheorem{theorem}{Theorem}
\newtheorem{lemma}{Lemma}
\newtheorem{proposition}{Proposition}
\newtheorem{corollary}{Corollary}
\newtheorem{definition}{Definition}
\newtheorem{example}{Example}
\title{ShadowBench algebra/L2/alg\_gen\_L2\_007}
\begin{document}
\maketitle
% Problem id: algebra/L2/alg_gen_L2_007
% Companion instructions: docs/instructions.md
% Candidate skeletons: docs/skeletons/
% Lean target: ShadowBench/Source/Main.lean

\begin{theorem}
Given polynomials $f,g,h$ in $k[x_1,\dots,x_n]$, where
$h=\gcd(f,g)$ means that $h$ divides both $f$ and $g$ and is divisible by every common divisor of
$f$ and $g$, prove that $h=\gcd(f,g)$ if and only if $\langle h\rangle$ is the smallest
principal ideal containing $\langle f,g\rangle$.
\end{theorem}
\end{document}
